\section{Mechanization and Trusted Base}
\label{sec:artifact}

The development contains 27,987 lines of Rocq, excluding the manuscript's
copied source excerpts.  Since a majority of the codebase was written with AI,
the line count measures the scope of the checked development rather than human
effort.
\Cref{fig:theory-map} separates the application-specific security argument
from the reusable probability and program-logic layers.

\begin{figure*}[t]
\centering
\begin{tikzpicture}[
  x=1cm,
  y=1cm,
  >={Latex[length=2mm]},
  dependency/.style={->,semithick,draw=black!65},
  module/.style={
    draw=black!65,
    rounded corners=2pt,
    align=center,
    text width=3.0cm,
    minimum height=1.35cm,
    inner sep=3pt,
    font=\footnotesize
  },
  interface/.style={module,fill=orange!10},
  application/.style={module,fill=green!10},
  reusable/.style={module,text width=3.35cm,fill=blue!8},
  result/.style={module,fill=violet!12,draw=black!80,very thick},
  lane/.style={font=\sffamily\footnotesize\bfseries,anchor=west}
]

% Application-specific path.
\node[lane] at (0,5.45) {APPLICATION-SPECIFIC SECURITY PROOF};

\node[interface] (interfaces) at (1.55,4.35) {%
  \textbf{Scheme and security interfaces}\\[0.2ex]
  \(S\), charts, correctness, and games};

\node[application] (games) at (5.15,4.35) {%
  \textbf{Noise-flooded construction and games}};

\node[application] (local) at (8.75,4.35) {%
  \textbf{One-call Gaussian bound and table invariant}};

\node[application] (adaptive) at (12.35,4.35) {%
  \textbf{Exact operation bridges and adaptive lossy hop}};

\node[result] (final) at (15.95,4.35) {%
  \textbf{Game reduction and checked theorem}\\[0.2ex]
  \(\beta_{\mathsf{CPA}}(\mathcal B)+
    \sqrt{qn}/(2\gamma)\)};

\draw[dependency] (interfaces) -- (games);
\draw[dependency] (games) -- (local);
\draw[dependency] (local) -- (adaptive);
\draw[dependency] (adaptive) -- (final);

% Reusable mathematical and program-logic path.
\node[lane] at (0,2.55) {REUSABLE FOUNDATIONS};

\node[reusable] (kl) at (1.85,1.35) {%
  \textbf{KL foundations}\\[0.2ex]
  Pinsker and the conditional-coordinate chain bound};

\node[reusable] (gaussian) at (6.15,1.35) {%
  \textbf{Discrete-Gaussian analysis}\\[0.2ex]
  Normalization, moments, and exact KL};

\node[reusable] (logic) at (10.45,1.35) {%
  \textbf{Semantic program logics}\\[0.2ex]
  Additive-error and Pythagorean rules};

\node[reusable] (compiler) at (14.75,1.35) {%
  \textbf{Selected-call compiler}\\[0.2ex]
  Adaptive replacement\\
  theorem};

\draw[dependency] (kl) -- (gaussian);
\draw[dependency] (gaussian) -- (logic);
\draw[dependency] (logic) -- (compiler);

% The two points at which the reusable layer discharges the application.
\draw[dependency] (gaussian.north) -- (local.south);
\draw[dependency] (logic.north) -- (local.south);
\draw[dependency] (compiler.north) -- (adaptive.south);

\end{tikzpicture}
\caption{Theory map of the Rocq development.  Each box is a conceptually
significant module or a small, tightly coupled module group; an arrow points
from a dependency to code that uses it.  The upper lane specializes the
reusable lower lane to noise flooding and culminates in the exported security
theorem.  Orange marks input interfaces, green the application-specific
security proof, blue reusable foundations, and purple the exported result.
Utility and adapter modules, repeated transitive imports, and proof-engineering
dependencies are omitted.}
\label{fig:theory-map}
\end{figure*}


The mechanization contains a few layers.  The scheme layer
defines encryption games and the noise-flooded transform.  The
probability layer develops facts about countable discrete distributions,
without reference to programs.  The logic layer lifts those facts to SSProve
code and proves the selected-call compiler sound.  Finally, the security layer
instantiates the generic logic with the decryption invariant, constructs the
IND-CPA reduction, and composes the exact and quantitative game transitions.
The top-level result is a functor over precisely the scheme, chart,
probability-one correctness, IND-CPA security, and flooding-parameter
interfaces described in \Cref{sec:construction}; probability lemmas and
program rules are proved dependencies.

\subsection{Formalization of judgments and rules}

We follow SSProve's convention and engineering framework: pure computation is
embedded in Rocq, while probabilistic and stateful effects are represented by
SSProve's free-monadic program syntax, packages, heaps, and linking operations
\cite{ssprove}.  The new logic does not extend Rocq with a primitive judgment
or a collection of trusted proof rules.
Judgments are formalized as ordinary Rocq propositions over the denotation \texttt{Pr\_code},
and inference rules are lemmas which prove such judgments.

For example, the exact definition of the additive-error judgment is shown
below.  Its types range over SSProve programs and heaps; \texttt{complete}
adds an explicit outcome for missing subdistribution mass.

\begingroup
\setmonofont{FreeMono}[Scale=MatchLowercase]
\inputminted[
  fontsize=\scriptsize,
  breaklines,
  linenos,
  firstnumber=57,
  firstline=57,
  lastline=69
]{coq}{rocq-excerpts/theories/ProgramLogics/Ae.v}
\endgroup

Thus, for every related pair of inputs and initial heaps, the proposition
requires a coupling of the completed output-and-heap distributions in which
the postcondition holds with probability at least \(1-\varepsilon\).
The Pythagorean judgment is encoded in the same style.  For each related input
pair it existentially quantifies two full transcript distributions, requires
their final marginals to equal the two completed program denotations, and
requires their supports to satisfy the postcondition.  Its predicate
\texttt{pythDist} additionally imposes full mass and the coordinatewise
conditional-KL bounds from \Cref{def:pyth-judgment}.
To place every result type in one transcript space, the Rocq definition
injectively encodes typed output/heap pairs in \texttt{nat * heap} and
represents \(\bot\) by \texttt{None}.

The rules in \Cref{fig:logic-rules} are proved implications between these
definitions.  For instance, the source-level statement of the final
KL-to-additive-error bridge is:

\begingroup
\setmonofont{FreeMono}[Scale=MatchLowercase]
\inputminted[
  fontsize=\scriptsize,
  breaklines,
  linenos,
  firstnumber=529,
  firstline=529,
  lastline=542
]{coq}{rocq-excerpts/theories/ProgramLogics/Pyth.v}
\endgroup

Here \texttt{pythagorean\_tv\_bound s} is
\(\sqrt{\sum_i s_i/2}\).  Its proof applies the verified probability theorem
to the transcript witnesses, projects to their final marginals, and invokes
the completed-distribution coupling theorem.  The sampling, sequencing,
mixed Hoare/Pythagorean, and compiler rules have the same status: each is a
Rocq lemma proved from the program semantics.  Consequently, a Rocq kernel check of
the security theorem also checks the soundness chain from the distribution
lemmas through the program judgments.

\subsection{Probability formalization}

The probability development works over MathComp's countable discrete
distributions, rather than restricting the Gaussian argument to finite
support.  It defines KL finiteness as absolute continuity together with
summability of the KL integrand, proves Pinsker's inequality, and proves the
conditional-coordinate chain bound used in
\Cref{thm:pyth-probability}.  It also constructs completion and maximal
couplings, so that a distribution-level total-variation bound yields the
program-level proposition above.

The distributions used here, notably the discrete Gaussians, have countably
infinite support.  Consequently, the KL chain rule involves infinite series.
A paper proof may leave their convergence implicit; the formalization must
establish it before decomposing transcript KL into conditional KL costs.  It
also treats histories of probability zero explicitly.  These results yield the
chain bound used in the square-root composition argument.

For the one-call bound, the development also proves that the discrete Gaussian
is normalized and has full support, that its normalizer is invariant under
integer shifts, and that its centered first moment is zero.  These facts yield
the exact equal-variance KL identity in \Cref{eq:dg-kl} and discharge the
Gaussian sampling rule.  Further theorem statements and proof structure
appear in \Cref{sec:verified-kl-analysis}.

\subsection{The checked theorem}

At the application boundary, the functor defines the concrete reduction
\(\cB_{\cA,q}\) and makes the assumed IND-CPA bound and the statistical loss
explicit:

\begingroup
\setmonofont{FreeMono}[Scale=MatchLowercase]
\inputminted[
  fontsize=\scriptsize,
  breaklines,
  linenos,
  firstnumber=530,
  firstline=530,
  lastline=533
]{coq}{rocq-excerpts/theories/Security/NoiseFloodingSecurity/Final.v}
\endgroup

After proving that the constructed reduction has the IND-CPA adversary
interface and deriving the completed-output inequality, the development
exports the following theorem.

\begingroup
\setmonofont{FreeMono}[Scale=MatchLowercase]
\inputminted[
  fontsize=\scriptsize,
  breaklines,
  linenos,
  firstnumber=1953,
  firstline=1953,
  lastline=1956
]{coq}{rocq-excerpts/theories/Security/NoiseFloodingSecurity/Final.v}
\endgroup

The premise says that \(A\) is a well-typed package with the required IND-CPAD
adversary interface.  Unfolding \texttt{security\_bound} and
\texttt{security\_loss} gives exactly \Cref{thm:main}; in particular, the
right-hand term is the assumed IND-CPA winning-probability bound for the
constructed reduction plus \(\sqrt{qn}/(2\gamma)\).  This source statement also
makes clear where the theorem boundary lies: scheme structure, charts,
probability-one correctness, IND-CPA security, and positivity of \(\gamma\)
enter through the enclosing functor, whereas reduction validity, the
discrete-Gaussian calculation, compiler correctness, and logic soundness have
already been proved.

\subsection{Reproducibility}

The complete artifact checks with Rocq 9.0.1, \texttt{rocq-ssprove}
development version at commit \texttt{c6d7d4bc3a}, MathComp Analysis 1.16.0,
and MathComp algebra tactics 1.2.7.  In a fresh switch on the authors' machine,
a clean build took about six minutes with four parallel jobs; this is a single
observation for reviewer planning, not a performance claim.  The manuscript's
Rocq listings are refreshed from the checked sources during its build, so the
displayed propositions and theorem statements do not silently diverge from
the artifact.

\subsection{Trusted computing base}
\label{sec:tcb}

The trusted base includes the Rocq kernel and standard library; MathComp,
MathComp Analysis, and their real-number/classical infrastructure; and
SSProve's program, heap, package, linking, and subdistribution semantics.
The scheme, chart, correctness, parameter, and IND-CPA interfaces from
\Cref{sec:construction} are theorem parameters and must be discharged by a
concrete instantiation.

The new KL, Gaussian, relational, compiler, and game-hop results contain no
local \texttt{Admitted}.  The installed dependency graph nevertheless reports
MathComp Analysis's
\texttt{realsum.\_\_admitted\_\_interchange\_psum}: current SSProve
distribution semantics are built on these MathComp definitions and reach the
upstream placeholder~\cite{ssprove}.  This is an inherited assumption, not an
unproven axiom introduced by our security development.

We independently discharge the mathematical obligation.  The artifact proves
the same statement as \texttt{interchange\_psum\_proved};
\texttt{Print Assumptions} for this replacement reports only the usual
extensionality and choice principles and not the admitted upstream lemma.  The
proof has since been merged to MathComp Analysis's development branch,
together with proofs of \texttt{psumB} and \texttt{interchange\_sup}, for
inclusion in the planned 1.17.0 release.  Until an upstream release replaces
the existing call sites, \texttt{Print Assumptions} for the top-level theorem
will continue to display the placeholder's name.

Besides library and classical assumptions, \texttt{Print Assumptions} for an
instantiated theorem reports the concrete scheme, chart, correctness, and
IND-CPA interfaces, together with positivity of the discrete-Gaussian width
multiplier \(\gamma\).  It does not report any program-logic rule, compiler
correctness theorem, or discrete-Gaussian KL identity: these are verified
inside the development.

Selected definitions of the scheme interface, games, transform,
and reduction remain in \Cref{app:formal-interface-excerpts}.  Together with
the source-level theorem above, they expose the small specification surface on
which a cryptographic audit should concentrate; the Rocq kernel checks that the
remaining proof terms connect that surface by the argument in
\Cref{sec:verified-reduction}.  As in SSProve and EasyCrypt, running time is not
part of the program semantics and must be inspected separately
\cite{easycrypt,ssprove}.  Here the
reduction runs the adversary once, forwards its oracle calls, and adds table
operations and discrete-Gaussian sampling.
